\documentclass{article}
\usepackage{amsmath,amssymb}
\usepackage[vocab={analyticnumbertheory}]{leanify}
\newtheorem{lemma}{Lemma}
\begin{document}

\section{Notation}

We write $\Lambda(n)$ for the von Mangoldt function at $n$.
We write $\vartheta(x)$ for the first Chebyshev function at $x$.
We write $\psi(x)$ for the second Chebyshev function at $x$.
We write $\pi(x;q,a)$ for the number of primes up to $x$ congruent to $a$ modulo $q$.
We write $\pi_2(x)$ for the number of twin prime pairs up to $x$.
We write $g_n$ for the $n$th prime gap.

\section{Values}

\begin{lemma}\label{mangoldt}
$\Lambda(8)=\log 2$, $\Lambda(6)=0$ and $\Lambda(7)=\log 7$.
\end{lemma}

\begin{lemma}\label{chebyshev}
$\vartheta(4)=\log 6$ and $\psi(4)=\log 12$.
\end{lemma}

\begin{lemma}\label{counting}
$\pi(20;4,1)=3$, $\pi_2(20)=4$ and $g_1=1$.
\end{lemma}

\begin{lemma}\label{comparison}
Let $x$ be a real number. Then $\vartheta(x)\le\psi(x)$.
\end{lemma}

\end{document}
